% !PRINTED 421
\weberpage{421}

Bilden wir aus diesen invarianten Formen des Tetra\"eders Functionen von $x_1:x_2$ und setzen $\vartheta = x_1:x_2$, so ergeben sich Functionen, die durch die linearen gebrochenen Tetra\"edersubstitutionen unge\"andert bleiben, und die daher Wurzeln von Resolventen 5\textsuperscript{ten} Grades der Ikosa\"edergleichung sind. Die Ikosa\"edergleichung selbst erh\"alt man in zwei Formen, wenn man
\weq{10}{H^3 = zf^5,\quad T^2 = z_1f^5,\quad z + z_1 = 1728}
setzt, und die Resolventen h\"angen rational von $z$ ab.

Bezeichnen wir f\"ur den Augenblick mit $\varphi(x_1, x_2)$ irgend eine absolut invariante Form des Tetra\"eders, so geht diese Form durch irgend eine der Substitutionen $Q\Theta$ in
\weq{11}{\varphi_\varepsilon = \varphi(\varepsilon^{-2}x_1,\ \varepsilon^2x_2)}
\"uber, und diese Function bleibt unge\"andert durch die Substitutionen der Gruppe $\Theta^{-1}Q\Theta$.

\gesperrt{Daraus ergiebt sich, dass jede symmetrische Function der f\"unf Formen $\varphi_\varepsilon$ eine Invariante der Ikosa\"edergruppe ist, und daher nach dem Satze \S.\,61 durch die Grundformen $f, H, T$ ausgedr\"uckt werden kann.}

Dieser Satz gestattet eine verh\"altnissm\"assig leichte Berechnung der Resolventen.

Die einfachsten Functionen, die wir als Wurzeln von Resolventen verwenden k\"onnen, sind
\weq{12}{r = \frac{t^2}{f},\quad u = \frac{f^2t}{T},\quad v = \frac{f\tilde\omega}{H}.}

Bemerken wir zun\"achst, dass die symmetrischen Grundfunctionen der f\"unf Functionen $t_\varepsilon$ Ikosa\"ederinvarianten der Grade $6, 12, 18, 30$ sind, so folgt aus den S\"atzen in \S.\,61, da die Gradzahlen $6, 18$ unter diesen Invarianten nicht vorkommen, eine Gleichung von der Form
\weq{13}{t^5 + aft^3 + bf^2t + cT = 0,}
worin $a, b, c$ Zahlenco\"efficienten sind. Um sie zu bestimmen, bezeichnen wir mit $S(\varphi)$ die Summe der f\"unf Functionen $\varphi_\varepsilon$ in (11) und mit $\Pi(\varphi)$ ihr Product, und erhalten aus den \gesperrt{Newton}'schen Formeln (Bd.\,I, \S.\,42):
\weq{14}{af = -\tfrac{1}{2}S(t^2),\quad bf^2 = \tfrac{1}{2}a^2f^2 - \tfrac{1}{4}S(t^4),\quad cT = -\Pi(t).}

In den Summen $S(\varphi)$ haben nur solche Glieder $x_1^h x_2^k$ einen von Null verschiedenen Co\"efficienten, in denen $h - k$ durch 5 theilbar ist. Man erh\"alt also $a, b, c$ durch Gleichsetzen der
%CONT>>
